\mychapter{0}{Résumé}


Les smartphones et montres connectées embarquent des capteurs inertiels qui
permettent de reconnaître ce que fait une personne : marcher, monter des
escaliers, s'asseoir, etc. En pratique, un modèle entraîné une fois en
laboratoire se dégrade dès qu'un nouvel utilisateur arrive ou que le contexte
change. C'est l'oubli catastrophique. Peu de systèmes, en outre, cherchent à
anticiper une transition avant qu'elle ne soit complète.

Nous proposons ici une architecture HAR qui apprend de façon continue et qui
peut aussi prédire l'activité suivante à partir de fenêtres partielles. Le
cœur du système est un Transformer partagé, un tampon de rejeu guidé par
l'incertitude de prédiction (UWR), et une tête LSTM d'anticipation. Sur
HAPT et WISDM, UWR gagne environ 16,8 points de macro-F1 face au finetuning
naïf en scénario user-incrémental, et dépasse EWC et iCaRL sur nos protocoles.
L'anticipation multi-classe atteint un macro-F1 de validation d'environ
0,60--0,64 (contre $\approx$0,04 pour une baseline majorité). Un module
binaire de risque postural (fall-risk) monte à environ 0,88.

Le code (\texttt{har\_project}) contient aussi un pré-entraînement SSL léger,
un adaptateur de domaine affine, une calibration de confiance en ligne et un
petit bandit de seuil d'alerte. Le backbone fait environ 531\,K paramètres ;
le tampon de rejeu est limité à 2\,000 fenêtres. Le transfert HAPT vers WISDM
reste difficile (F1 $\approx$0,46 en zéro-shot). Nous n'avons pas de corpus
cuisine pour le scénario « préparation de repas » évoqué dans la fiche de
stage. Les annexes détaillent les données, les hyperparamètres et le code
pour reproduire les expériences.

\bigskip
\noindent\textbf{Mots clés :} Reconnaissance d'activités humaines, apprentissage
continu, capteurs inertiels, Transformer, Uncertainty-Weighted Replay, anticipation
d'activités.

\clearpage

\mychapter{0}{Abstract}


Inertial sensors in phones and watches make it possible to recognise daily
activities, but a model trained once on laboratory data usually degrades when
new users or contexts appear. Few systems try to anticipate a transition
before it is complete.

This thesis describes a continual HAR system built around a shared Transformer,
an uncertainty-weighted replay buffer (UWR), and an LSTM head that predicts the
next activity from partial windows. On HAPT and WISDM, UWR improves macro-F1 by
about 16.8 points over naive fine-tuning in the user-incremental setting and
outperforms EWC and iCaRL in our protocols. Next-activity anticipation reaches
validation macro-F1 around 0.60--0.64 (majority baseline $\approx$0.04). A
binary postural fall-risk head reaches about 0.88. Lighter modules cover SSL
pretraining, a domain adapter, online confidence calibration and a threshold
bandit. Cross-dataset transfer remains hard (HAPT$\rightarrow$WISDM F1
$\approx$0.46 zero-shot).

\bigskip
\noindent\textbf{Keywords:} Human activity recognition, continual learning, inertial
sensors, Transformer, Uncertainty-Weighted Replay, activity anticipation.

\clearpage

\chapter*{\hfill \textarabic{ملخص}}

\begin{otherlanguage}{arabic}
\noindent
\textarabic{%
يمكن لحساسات القصور الذاتي في الهواتف والساعات التعرف على الأنشطة اليومية،
لكن النموذج المدرّب مرة واحدة على بيانات مخبرية يتراجع أمام مستخدمين أو سياقات
جديدة، وقليل من الأنظمة يتوقع الانتقال قبل اكتماله.}

\bigskip
\noindent
\textarabic{%
يصف هذا المذكّر نظام HAR مستمراً يعتمد على Transformer مشترك، وإعادة تشغيل
مرجّحة بعدم اليقين (UWR)، ورأس LSTM لتوقع النشاط التالي. التجارب على HAPT
وWISDM تُظهر تحسناً واضحاً مقارنة بالضبط البسيط وEWC وiCaRL.}

\bigskip
\noindent\textbf{\textarabic{الكلمات المفتاحية:}}
\textarabic{%
التعرف على الأنشطة البشرية، التعلم المستمر، حساسات القصور الذاتي،
Transformer، UWR، توقع النشاط.}
\end{otherlanguage}

\addcontentsline{toc}{chapter}{ملخص}
